\documentclass[a4paper,12pt]{article}
\usepackage[T2A]{fontenc}
\usepackage[utf8]{inputenc}
\usepackage[russian]{babel}
\usepackage{amsmath,amssymb,mathtools}
\renewcommand{\sfdefault}{cmss}
\renewcommand{\familydefault}{\sfdefault}
\usepackage{icomma}
\usepackage{geometry}
\geometry{left=20mm,right=20mm,top=20mm,bottom=21mm}
\usepackage{microtype}
\usepackage{tikz}
\usetikzlibrary{calc,arrows.meta,positioning,shapes.geometric}
\usepackage{pgfplots}
\pgfplotsset{compat=1.18}
\usepackage{tabularx,booktabs}
\usepackage{enumitem}
\usepackage{fancyhdr}
\usepackage{setspace}
\usepackage{xcolor}
\usepackage[hidelinks]{hyperref}
\definecolor{accent}{HTML}{008080}
\definecolor{darkaccent}{HTML}{004D4D}
\definecolor{textgray}{HTML}{30373A}
\definecolor{lightgray}{HTML}{657578}
\definecolor{softline}{HTML}{CFDDDA}
\color{textgray}
\onehalfspacing
\setlength{\parindent}{0pt}
\setlength{\parskip}{5pt}
\setlist[enumerate]{leftmargin=*,itemsep=7pt,topsep=5pt}
\setlist[itemize]{leftmargin=*,itemsep=4pt,topsep=3pt}
\emergencystretch=2em
\raggedbottom
\renewcommand{\today}{16 октября 2026 г.}
\pagestyle{fancy}
\fancyhf{}
\fancyhead[L]{\footnotesize\color{lightgray} СТЕПЕНИ \enspace / \enspace ПРАКТИКА}
\fancyhead[R]{\footnotesize\color{lightgray} 16.10.2026}
\fancyfoot[C]{\footnotesize\color{lightgray}\thepage}
\renewcommand{\headrulewidth}{0.3pt}
\renewcommand{\headrule}{\hbox to\headwidth{\color{softline}\leaders\hrule height \headrulewidth\hfill}}
\setlength{\headheight}{23pt}
\newcommand{\sect}[1]{\vspace{5pt}{\large\bfseries\color{darkaccent}#1}\par\vspace{3pt}{\color{softline}\hrule height .5pt}\vspace{6pt}}
\newcommand{\mini}[1]{{\bfseries\color{darkaccent}#1}\enspace}
\newcommand{\task}[2]{\item \textbf{#1}\par\vspace{2pt}#2}
\begin{document}
\thispagestyle{empty}
\begin{center}
\vspace*{34mm}
{\small\color{accent}\bfseries МАТЕМАТИКА \quad / \quad ПОВТОРЕНИЕ К ЕГЭ}\par
\vspace{18mm}
{\fontsize{32}{38}\selectfont\bfseries\color{darkaccent} Чек-лист}\par
\vspace{5mm}
{\LARGE Степени и показательная функция}\par
\vspace{3mm}
{\large Вычисления, прикладные задачи, графики}\par
\vspace{18mm}
{\color{softline}\rule{85mm}{0.7pt}}\par
\vspace{13mm}
{\large\bfseries Денис Ермолаев}\par
\vspace{4mm}
{\normalsize\today}\par
\vfill
{\normalsize Лёвин Артём Александрович}\par
{\small\color{lightgray} эксклюзивно для Дениса Ермолаева}\par
\vspace{21mm}
\begin{tikzpicture}[x=1cm,y=1cm]
\draw[accent!65,line width=1.15pt] (-6,0.15) .. controls (-3.5,0.95) and (-1.5,-0.25) .. (0.9,0.36) .. controls (3,0.9) and (4.6,-0.2) .. (6.1,0.3);
\end{tikzpicture}
\end{center}
\clearpage

\sect{01 \quad Опорные формулы}
Для положительных оснований $a,b>0$ и действительных показателей $m,n$ верны:
\renewcommand{\arraystretch}{1.47}
\begin{tabularx}{\linewidth}{@{}p{0.45\linewidth}X@{}}
\toprule
\textbf{Действие} & \textbf{Свойство}\\
\midrule
Умножение одинаковых оснований & $a^m\cdot a^n=a^{m+n}$\\
Деление одинаковых оснований & $\dfrac{a^m}{a^n}=a^{m-n}$\\
Степень произведения & $(ab)^n=a^nb^n$\\
Степень частного & $\left(\dfrac ab\right)^n=\dfrac{a^n}{b^n}$\\
Степень степени & $(a^m)^n=a^{mn}$\\
Нулевая и отрицательная степени & $a^0=1$,\quad $a^{-n}=\dfrac1{a^n}$\\
Корень и дробный показатель & $\sqrt[k]{a^m}=a^{m/k}$, $k\in\mathbb N$, $k\ge2$\\
\bottomrule
\end{tabularx}

\mini{В обратную сторону тоже.} Например, $a^{m+n}=a^m a^n$ и $a^m b^m=(ab)^m$. Иногда именно обратное применение сокращает запись.

\mini{Смена числителя и знаменателя.} $\left(\dfrac ab\right)^n=\left(\dfrac ba\right)^{-n}$ при $a,b>0$.

\mini{Об ограничениях.} При произвольных действительных показателях пользуйтесь формулами с положительными основаниями. Для целых показателей основания могут быть отрицательными; знаменатель и основание отрицательной степени должны быть ненулевыми.

\sect{02 \quad Главный приём вычислений}
\textbf{Составные числа, корни и дроби приводим к степенному виду, затем к общему основанию.}

\begin{enumerate}
\item \mini{Составное число.} $25=5^2$,\quad $81=3^4$,\quad $32=2^5$.
\item \mini{Корень.} $\sqrt[3]{2}=2^{1/3}$,\quad $\sqrt{2^3}=2^{3/2}$.
\item \mini{Дробь.} $0{,}5=\dfrac12=2^{-1}$,\quad $\dfrac1{\sqrt[3]{81}}=3^{-4/3}$.
\item \mini{Общее основание.} Сложите показатели при умножении, вычтите при делении, перемножьте при возведении степени в степень.
\end{enumerate}
\mini{Пример.} \(5^{0{,}36}\cdot 25^{0{,}32}=5^{0{,}36}\cdot(5^2)^{0{,}32}=5^{0{,}36+0{,}64}=\boxed{5}.\)

\mini{Скобки («ВУПС»).} Вычитаемое, умножаемое и подставляемое выражения заключайте в скобки. Например, $a^{x-(-y)}=a^{x+y}$.

\clearpage
\sect{03 \quad Как исчезают корни и дробные показатели}
При приведении степеней к одному основанию иррациональные показатели могут сократиться.

\mini{Пример.}
\[
\frac{(0{,}5)^{\sqrt{10}-1}}{2^{-\sqrt{10}}}
=\frac{(2^{-1})^{\sqrt{10}-1}}{2^{-\sqrt{10}}}
=2^{-\sqrt{10}+1-(-\sqrt{10})}=\boxed{2}.
\]

\mini{Дробные показатели.}
\[
\left(\frac{2^{1/3}\,\sqrt[4]{2}}{\sqrt[12]{2}}\right)^2
=\left(2^{1/3+1/4-1/12}\right)^2
=2^{(1/2)\cdot 2}=\boxed{2}.
\]

\sect{04 \quad Прикладные задачи со степенями}
\mini{1. Формула газа: найти величину.} Пусть $pV^k=C$, где $p,V,C>0$, $k\ne0$. Чтобы выразить объём, возведите обе части равенства $V^k=C/p$ в степень $1/k$:
\[
\boxed{\;V=\left(\frac Cp\right)^{1/k}\;}
\]
В примере $pV^{5/3}=10^5$, $p=3{,}2\cdot10^6$:
\[
V^{5/3}=\frac{10^5}{32\cdot10^5}=2^{-5},\qquad
V=(2^{-5})^{3/5}=\boxed{\tfrac18=0{,}125}.
\]
\mini{Смысл дробного показателя.} Чтобы получить $V$ из $V^{5/3}$, обе части возводим в степень $3/5$: показатели перемножаются и дают единицу.

\mini{2. Радиоактивный распад.} При $m=m_0\cdot2^{-t/T}$ ($m_0>0$, $T>0$) подставьте исходную и конечную массу, а затем приведите обе части к основанию $2$.
\[
\frac{5}{40}=2^{-t/10}\quad\Longrightarrow\quad
2^{-3}=2^{-t/10}\quad\Longrightarrow\quad \boxed{t=30\ \text{мин}}.
\]
Равенство $2^u=2^v$ позволяет заключить $u=v$, поскольку основание $2>0$ и $2\ne1$.

\clearpage
\sect{05 \quad Сравнение двух состояний}
Если $pV^{\alpha}=\text{const}$, для двух состояний имеем
\[
\boxed{\;p_1V_1^{\alpha}=p_2V_2^{\alpha}\;}.
\]
При уменьшении объёма в 2 раза $V_2=V_1/2$, а при увеличении давления в 4 раза $p_2=4p_1$. Тогда при $p_1,V_1>0$:
\[
p_1V_1^\alpha=(4p_1)\left(\frac{V_1}{2}\right)^\alpha
\quad\Longrightarrow\quad 2^\alpha=4=2^2
\quad\Longrightarrow\quad \boxed{\alpha=2}.
\]
\textbf{Важно:} «уменьшился в $q$ раз» означает деление на $q$, а «увеличился в $q$ раз» --- умножение на $q$. Индексы $1$ и $2$ обозначают состояния до и после изменения.

\sect{06 \quad Метод опорных (узловых) точек}
Для показательной функции $f(x)=a^x+b$ ($a>0$, $a\ne1$) нужно найти два неизвестных параметра $a$ и $b$. По графику удобно выбирать \textbf{две точки с точно читаемыми координатами}, особенно при $x=0$ и $x=1$.
\[
f(0)=1+b,\qquad f(1)=a+b.
\]
\mini{Пример.} По точкам $(0;-2)$ и $(1;-1)$ получаем
\[
\begin{cases}
-2=1+b,\\[-1mm]
-1=a+b,
\end{cases}\qquad b=-3,\quad a=2.
\]
Значит, $f(x)=2^x-3$, откуда $f(6)=2^6-3=\boxed{61}$.

\begin{center}
\begin{tikzpicture}
\begin{axis}[width=11.5cm,height=6.5cm,xmin=-2,xmax=2.35,ymin=-3.5,ymax=3.6,
axis lines=middle,axis line style={->,black!65},xtick={-2,-1,0,1,2},ytick={-3,-2,-1,0,1,2,3},grid=both,
minor tick num=0,grid style={softline,thin},xlabel={$x$},ylabel={$y$},xlabel style={at={(axis description cs:1,.45)},anchor=south},ylabel style={at={(axis description cs:.49,1)},anchor=west},
clip=false, samples=201,tick label style={font=\small}]
\addplot[very thick,accent,domain=-2:2.2] {2^x-3};
\addplot[only marks,mark=*,mark size=2.2pt,darkaccent] coordinates {(0,-2)(1,-1)};
\node[anchor=east,font=\small,fill=white,inner sep=2pt] at (axis cs:-.1,-2) {$(0;-2)$};
\node[anchor=south west,font=\small,fill=white,inner sep=2pt] at (axis cs:1,-1) {$(1;-1)$};
\end{axis}
\end{tikzpicture}
\end{center}
\mini{Самопроверка.} Для $a>1$ график возрастает, для $0<a<1$ --- убывает. Горизонтальная асимптота: $y=b$. Обязательно сверяйте найденную формулу с видом графика.

\clearpage
\thispagestyle{fancy}
\begin{center}
{\Large\bfseries\color{darkaccent} Алгоритм: вычисление выражений со степенями}\par
\vspace{9mm}
\begin{tikzpicture}[>=Stealth,thick,node distance=13mm,
start/.style={ellipse,draw=darkaccent,fill=accent!5,align=center,text width=57mm,minimum height=12mm,inner sep=5pt},
action/.style={rounded corners=3pt,rectangle,draw=softline,fill=white,align=center,text width=75mm,minimum height=15mm,inner sep=7pt},
decision/.style={diamond,aspect=2.3,draw=darkaccent,align=center,text width=53mm,inner sep=4pt},
arrow/.style={-Stealth,color=darkaccent,line width=.95pt}]
\node[start] (st) {Прочитайте выражение};
\node[decision,below=19mm of st] (dec) {Есть составные основания,\\корни или дроби?};
\node[action,right=20mm of dec,text width=52mm] (norm) {Представьте их степенями\\и приведите основания\\к удобному виду};
\node[action,below=24mm of dec] (rules) {Примените правила степеней:\\умножение, деление,\\степень степени};
\node[action,below=14mm of rules] (simple) {Приведите показатели:\\учтите знаки, скобки\\и дробную арифметику};
\node[start,below=15mm of simple] (done) {Вычислите результат\\и проверьте ограничения};
\draw[arrow] (st) -- (dec);
\draw[arrow] (dec.south) -- node[left,font=\small]{нет} (rules.north);
\draw[arrow] (dec.east) -- node[above,font=\small]{да} (norm.west);
\draw[arrow] (norm.south) |- ([yshift=0mm]rules.east);
\draw[arrow] (rules) -- (simple);
\draw[arrow] (simple) -- (done);
\end{tikzpicture}
\end{center}
\clearpage

\sect{07 \quad Тренировочный блок \hfill Путь А}
{\small\color{lightgray}10 заданий. Упражнения расположены по нарастающей сложности. Покажите основные преобразования и отдельно запишите ответ.}\par
\begin{enumerate}[label=\textbf{\arabic*.},leftmargin=8mm,itemsep=11pt]
\task{Вычислите:}{\[7^{0{,}4}\cdot49^{0{,}3}.\]}
\task{Вычислите:}{\[\left(\frac{3^{1/4}\cdot\sqrt[6]{3}}{\sqrt[12]{3}}\right)^{3}.\]}
\task{Вычислите:}{\[\frac{(0{,}25)^{\sqrt{3}-2}}{4^{-\sqrt{3}}}.\]}
\task{Вычислите:}{\[\frac{9^{\sqrt{7}-1}}{3^{2\sqrt{7}-4}}.\]}
\task{Вычислите:}{\[\left(\frac{1}{\sqrt[3]{64}}\right)^{-3/2}.\]}
\task{Газ и степень.}{Величины связаны формулой $pV^{5/3}=10^7$. Найдите $V$, если $p=3{,}125\cdot10^5$.}
\task{Радиоактивный распад.}{Масса вещества изменяется по закону $m=m_0\cdot2^{-t/T}$, где $t$ измеряется в минутах. При $m_0=96$ г и $T=8$ мин найдите время, через которое масса станет равна $12$ г.}
\task{Два состояния газа.}{В процессе $pV^\alpha=\mathrm{const}$ объём уменьшился в $3$ раза, а давление увеличилось в $9$ раз. Найдите $\alpha$, считая давление и объём положительными.}
\end{enumerate}
\clearpage
\sect{07 \quad Тренировочный блок \hfill Продолжение}
\begin{enumerate}[label=\textbf{\arabic*.},leftmargin=8mm,start=9,itemsep=14pt]
\task{Показательная функция (возрастание).}{На графике изображена функция $f(x)=a^x+b$. Используя отмеченные точки, найдите $f(3)$.}
\end{enumerate}
\begin{center}
\begin{tikzpicture}
\begin{axis}[width=12.2cm,height=6.7cm,xmin=-1.5,xmax=2.1,ymin=-5.4,ymax=5,
axis lines=middle,axis line style={->,black!65},xtick={-1,0,1,2},ytick={-5,-4,-3,-2,-1,0,1,2,3,4},
grid=both,grid style={softline,thin},minor tick num=0,xlabel={$x$},ylabel={$y$},
clip=false,tick label style={font=\small},xlabel style={at={(axis description cs:1,.51)},anchor=south},ylabel style={at={(axis description cs:.43,1)},anchor=west}]
\addplot[accent,very thick,domain=-1.5:1.96,samples=180] {3^x-5};
\addplot[only marks,mark=*,mark size=2.6pt,darkaccent] coordinates {(0,-4)(1,-2)};
\node[anchor=north east,font=\small,fill=white,inner sep=2pt] at (axis cs:0,-4) {$(0;-4)$};
\node[anchor=south east,font=\small,fill=white,inner sep=2pt] at (axis cs:1,-2) {$(1;-2)$};
\end{axis}
\end{tikzpicture}
\end{center}
\vspace{2mm}
\begin{enumerate}[label=\textbf{\arabic*.},leftmargin=8mm,start=10,itemsep=12pt]
\task{Показательная функция (убывание).}{На графике изображена функция $g(x)=a^x+b$. По отмеченным точкам определите $g(-3)$.}
\end{enumerate}
\begin{center}
\begin{tikzpicture}
\begin{axis}[width=12.2cm,height=6.3cm,xmin=-2,xmax=2.35,ymin=-4.6,ymax=0.4,
axis lines=middle,axis line style={->,black!65},xtick={-2,-1,0,1,2},ytick={-4,-3,-2,-1,0},
grid=both,grid style={softline,thin},minor tick num=0,xlabel={$x$},ylabel={$y$},
clip=false,tick label style={font=\small},xlabel style={at={(axis description cs:1,.94)},anchor=south},ylabel style={at={(axis description cs:.47,1)},anchor=west}]
\addplot[accent,very thick,domain=-2:2.2,samples=180] {0.5^x-4};
\addplot[only marks,mark=*,mark size=2.6pt,darkaccent] coordinates {(-1,-2)(0,-3)};
\node[anchor=south east,font=\small,fill=white,inner sep=2pt] at (axis cs:-1,-2) {$(-1;-2)$};
\node[anchor=south west,font=\small,fill=white,inner sep=2pt] at (axis cs:0,-3) {$(0;-3)$};
\end{axis}
\end{tikzpicture}
\end{center}

\clearpage
\sect{08 \quad Ответы для самопроверки}
\vspace{5mm}
\renewcommand{\arraystretch}{1.65}
\begin{center}
\begin{tabular}{@{}c p{57mm} c@{}}
\toprule
\textbf{№} & \textbf{Проверяемая идея} & \textbf{Ответ}\\
\midrule
1 & Составное основание & $7$\\
2 & Корни и дробные степени & $3$\\
3 & Дробь и отрицательный показатель & $16$\\
4 & Приведение к общему основанию & $9$\\
5 & Отрицательная дробная степень & $8$\\
6 & Степень в формуле газа & $8$\\
7 & Закон распада & $24$ мин\\
8 & Сравнение двух состояний & $\alpha=2$\\
9 & Опорные точки возрастающего графика & $22$\\
10 & Опорные точки убывающего графика & $4$\\
\bottomrule
\end{tabular}
\end{center}
\vspace{9mm}
\sect{Как проверить своё решение}
\begin{itemize}
\item В № 1--5 основание после преобразований должно быть положительным, а дробные показатели --- корректно приведёнными.
\item В № 6--8 подставьте найденное значение в исходную формулу и проверьте отношения величин.
\item В № 9--10 найденные параметры $a,b$ должны удовлетворять координатам обеих отмеченных точек.
\end{itemize}
\vfill
{\small\color{lightgray}Лёвин Артём Александрович \quad · \quad эксклюзивно для Дениса Ермолаева}
\end{document}
